% front/01_cover.tex
% 모든 페이지와 동일 geometry(4/4/5cm) 사용

\thispagestyle{empty}

\begingroup
\setlength{\parskip}{0pt}
\setlength{\parindent}{0pt}

% --------------------------------
% 0) 박사 학위 논문  – 상단에서 0.5cm
% (top=4cm 이 이미 적용되어 있으니, 텍스트 블록 위에서 0.5cm만 더 내려감)
% --------------------------------
\vspace*{0.5cm}
{\fontsize{14}{20}\selectfont 박사 학위 논문}

% --------------------------------
% 1) 국문 제목 – 상단에서 3.5cm
% → 이미 0.5cm 내려왔으니, 3.5 - 0.5 = 3.0cm 더 내려주면 됨
% --------------------------------
\vspace*{3.0cm}

\begin{center}
  \begin{minipage}{13cm}  % 한글/영문 제목 폭 = 13cm
    \centering
    {\fontsize{18}{18}\selectfont\bfseries
    관측성과 불확실성 기반 동적\par
    복구목표시간(DRTO) 프레임워크 연구\par}
  \end{minipage}
\end{center}

% --------------------------------
% 2) 영문 제목 – 한글 제목 블록과 1cm 간격
% --------------------------------
\vspace*{1.0cm}

\begin{center}
  \begin{minipage}{13cm}
    \centering
    {\fontsize{18}{18}\selectfont\bfseries
      A Dynamic RTO Framework\par
      Based on Observability and Uncertainty\par}
  \end{minipage}
\end{center}

% --------------------------------
% 3) 아래 정보 (날짜 ~ 이름)
%    간격: 날짜-숭실대 3.5cm, 숭실대-학과 1.5cm, 학과-이름 1.5cm
%    "박스 하단 → 이름 1.5cm"는 보고 눈으로 미세조정하는 쪽이 현실적이라
%    아래 마지막 \vspace* 값을 조금씩 조절해 쓰시면 됩니다.
% --------------------------------

% 영문 제목 아래에서 날짜까지 여백 (눈으로 맞춰볼 값, 대략 3.0~3.5cm 정도에서 시작)
\vspace*{3.0cm}

\begin{center}
  {\fontsize{14}{20}\selectfont 2026년 6월}\par

  \vspace{2.0cm} % 날짜 → 숭실대 대학원 (3.5cm)
  {\fontsize{16}{22}\selectfont\bfseries 숭실대학교 대학원}\par

  \vspace{1.0cm} % 숭실대 대학원 → 학과 (1.5cm)
  {\fontsize{14}{20}\selectfont 재난안전관리학과}\par

  \vspace{1.0cm} % 학과 → 이름 (1.5cm)
  {\fontsize{14}{20}\selectfont 이\ \ \ 창\ \ \ 호}\par

  % 박스 하단에서 이름까지 1.5cm 목표
  % 부족하거나 넘치면 여기 값을 ±0.5cm 정도씩 조절해 보세요.
  \vspace*{1.5cm}
\end{center}

\endgroup

\clearpage
